\documentclass[11pt,a4paper]{ctexart}
\usepackage[margin=2.35cm]{geometry}
\usepackage{amsmath,amssymb,amsthm,mathtools,enumitem,needspace}
\usepackage[hidelinks]{hyperref}
\setlength{\parskip}{0.35em}
\setlength{\emergencystretch}{2em}
\setlist[enumerate]{leftmargin=2em,itemsep=0.2em,topsep=0.3em}
\allowdisplaybreaks[2]
\newcommand{\R}{\mathbb R}
\newcommand{\C}{\mathbb C}
\newcommand{\Z}{\mathbb Z}
\newcommand{\N}{\mathbb N}
\newcommand{\dd}{\,\mathrm d}
\newcommand{\ii}{\mathrm i}
\newcommand{\inner}[2]{\left\langle#1,#2\right\rangle}
\newcommand{\src}[1]{\par\noindent{\small\textbf{出处与处理：}#1}\par}
\newcommand{\idea}[1]{\par\noindent\textbf{证明思路：}#1\par}
\newcommand{\ex}[3]{\Needspace{6\baselineskip}\subsection*{#1\quad #2}\label{ex:#1}\src{#3}}
\theoremstyle{definition}
\newtheorem{result}{结论}[section]
\newcommand{\sources}{
\noindent\textbf{范围与来源}\quad 仅含 Fourier 级数，不含 Fourier 变换、含参积分、偏微分方程及数论应用。
\begin{enumerate}[label=\arabic*.]
\item \textbf{课本：}《高等微积分学》（文件日期2023.2.6），第5.1--5.3节及“习题五”中级数部分；原书第105--122、126--129页，PDF第106--123、127--130页。第122页仅取第5.3节结尾，不含该页开始的第5.4节。
\item \textbf{史：}常庚哲、史济怀《数学分析教程》第三版下册，第17.1--17.4节，原书及该文件PDF第303--339页。作业照片所用版本编号为第21章；本材料新增选题按实际核对的第三版第17章引用，不混用题号。
\item \textbf{Stein：}Stein--Shakarchi, \emph{Fourier Analysis: An Introduction}，从第2、3章中筛选贴近系数计算、核估计、收敛与Parseval的Exercises；Problems只取相关的简单题或前几问，具体题号与页码逐题注明。抽象函数空间、测度理论及独立的Abel、Tauber求和题不纳入本册。
\item \textbf{作业与笔记：}第三次作业（10.15）、第四次作业（10.22）及其教材照片；已录入的《傅里叶级数笔记》。作业与教材重合的题合并注明。课堂笔记中含参积分、Laplace、Watson等内容均不在本次范围。
\end{enumerate}
原书页码与PDF页码均从实际文件核对。翻译题、节选题、改编题和新增题分别注明；题面不含解题提示。来源文件的实际位置及后续修改入口见同目录说明。}
